\section{结论}
\label{sec:conclusion}
在明确的操作性类别映射下，原始掩膜审计表明，主要 LoveDA 源缺少精确草本监督。在指定 CLIP 探针和更新预算下，加入所考察的 FLAIR 源混合能够恢复 OpenEarthMap rangeland 性能，但同时降低源支持类别性能。历史原生 U-Net 记录在同一目标上重复了这一方向；另一目标上的评估和反向源混合，则在各自明确的依赖关系与来源限制内扩展了描述性证据。这些结果为已有标签空间和负迁移研究补充了具体的审计与受控经验分析。正例源标签覆盖、缺失类别恢复和源支持类别保持应当共同报告。现有实验既不能单独识别标签的因果作用，也没有证明已经解决这一权衡。
